\scanpage{11}
\mspage{7\deleted{\unsure{6|5}}\handwritten{7}}

\deleted{est seul à connaître et à manier, une langue qui d'instant en instant se forme et se transforme sous ses mains de magicien. Invisible, insaisissable, en\handwritten{f}ant espiègle \handwritten{et} \deleted{sage débonnaire} vieux sage énigmatique \deleted{\unsure{et}} \deleted{débonnaire} - c'est \uline{lui}, le Maître du verbe vivant, de la langue sans paroles, \deleted{où \unsure{dans}}\inserted{la \mbox{l\,a\,n\,g\,u\,e\,.\ -\ m\,è\,r\,e} où} les mots innombrables des innombrables langues-mots viennent \inserted{depuis la nuit des âges} puiser \deleted{leur sens,} leur \deleted{\unsure{force}}\handwritten{\inserted{vigueur}} et leur \handwritten{sens}.}

\margin{\handwritten{\unsure{espacer|espaces}}}

C'est une chose étrange que la langue-mère, la langue commune à tous les êtres humains (voire\handwritten{\inserted{même}}, à tous les êtres \ldots), \handwritten{\inserted{qu'elle}} soit à tel point inapte, pourrait-il sembler, comme mode de \uline{communication} \uline{d'un} \uline{sens}, d'\deleted{\unsure{u}}un être à un autre. Ce n'est pas, il est vrai, la langue du patron, et si elle est moyen de communication, ce n'est sûrement pas d'un patron à un autre.

Nous pouvons, bien sûr, la ``traduire'' dans le langage des mots, \deleted{\unsure{et}} \handwritten{C}'est bien ce que je ne manque pas de faire avec \deleted{le}\handwritten{un} \deleted{plus grand} soin\handwritten{\inserted{infini}}, chaque fois (ou presque \ldots) qu'un rêve \deleted{semble} m'interpelle\deleted{\unsure{r}} avec une force exceptionnelle. Une telle ``traduction'' est un moyen d'écoute, tel un stét\handwritten{\inserted{h}}oscope pour des oreilles un tantinet sourdes ! \deleted{Visiblement,} C'est bien utile et même indispensable, vu le cas\handwritten{.} \handwritten{M}ais il est \deleted{\unsure{s}}sûr aussi que la langue-images n'\deleted{\unsure{x}}est pas \deleted{plus} faite pour être traduite en mots, \handwritten{\inserted{pas plus}} que la langue-\deleted{\unsure{paroles}}\handwritten{mots} n'est faite pour le stét\handwritten{\inserted{h}}oscope. \deleted{\unsure{La}}\handwritten{\inserted{Et Toute}} traduction est aussi différente de la parabole originelle, qu'une description d'un feu, d'une eau ou d'une scène vécue, sont différents de la chose décrite.

Si j'essaye d'évoquer quelque moyen de communication directe par la langue-images, ou par quelque \handwritten{\inserted{autre}} ``langue'' qui s'en rapproche tant soit peu, je ne vois guère (faute d'imagination sans doute) que la peinture ou la sculpture, et surtout la \mbox{d\,a\,n\,s\,e} - la langue-mouvement par excellence, la langue du \mbox{c\,o\,r\,p\,s}.

Il est vrai que le corps a mille autres\deleted{\unsure{f}} façons encore de s'exprimer, de ``parler'' - \handwritten{par} un langage plus éloquent parfois que le langage des mots ou même la danse. Il y a le jeu d'amour, certes, ce Jeu des jeux, le jeu \deleted{de}\handwritten{\inserted{que joue}} la Terre \deleted{\unsure{avec|et} de}\handwritten{\inserted{avec}} ses amants innombrables, \handwritten{\inserted{avec}} le soleil, \handwritten{avec} le ciel \deleted{et}\handwritten{\inserted{avec}} la pluie, et \handwritten{\inserted{avec}} \deleted{toutes}\handwritten{\inserted{chacune de}} ses créatures \ldots{} Et il y a la langue des yeux\handwritten{\inserted{aussi}}, comme il y a celle des mains (sans compter celle des pieds\deleted{sous la table}, me souffle un petit diable blagueur \ldots),

\handwritten{D}écidément la liste \handwritten{\inserted{déjà}} s'étoffe! Et je pensais aussi, et surtout, au langage (que la Science baptise de ``psychosomatique'') par lequel le corps exprime et par là, peut-être, ``compense'' \deleted{en quelque sorte}\handwritten{\inserted{à sa façon}}, la violence subie par lui, ou par les couches profondes de la psyché qui s'enracinent en lui, aux mains de forces inquiètes et impitoyables, battant \handwritten{\inserted{souvent}} pavillon ``Esprit''. Et peut-être est-il vrai qu'il n'y a de ``sens'' bel et bien exprimé, \handwritten{\inserted{par qui que ce soit et}} dans quelque langue que ce soit, qui ne soit aussi bel et bien ``entendu'' \ldots

% archive sheet number: 10
